
T-cells play a key role in the adaptive immune system, contributing to immune surveillance and
response to pathogens. T-cells recognize pathogen-derived epitopes in the form of short protein
fragments (peptides) displayed on specific molecules known as Major Histocompatibility Complexes
(MHCs). The binding between a T-cell receptor (TCR) and a peptide-MHC (pMHC) is essential to mediate
a protective T-cell response to an infection. To defend the host against a multitude of pathogens, a
highly diverse TCR repertoire is generated through a somatic VDJ recombination process, and selected
to target infected cells with remarkable sensitivity and specificity.

Predicting the TCR-pMHC binding specificity is an important step in characterizing the efficacy of
the immune system in countering different pathogens. Such understanding would aid with disease
diagnoses and development of diagnostic tests for early detection of autoimmune
diseases~\cite{Attaf2015-tw}. A predictive model for TCR-pMHC specificity can be used to design
engineered TCRs that recognize cancer antigens~\cite{Liu2025-hl}, enhancing the effectiveness of
adoptive cell transfer and {Chimeric Antigen Receptor (CAR) T-cell}
therapies~\cite{Leidner2022-ca,Mirazee2024-bo}, as well as the design of soluble TCRs and bispecific
T-cell engager therapeutics~\cite{Hsiue2021-ag}. Moreover, it can enable the design of optimized
antigens for vaccines to elicit robust T-cell responses against specific pathogens or
tumors~\cite{Sahin2018-fv,Blass2021-jd,Rojas2023-nx}.

One bottleneck in deciphering the TCR-pMHC code is the lack of well-curated datasets.
High-throughput immune repertoire sequencing, combined with comparative analyses of hosts facing
similar immune challenges, and the development of functional assays (e.g., {pMHC-tetramer
staining}~\cite{Zhang2018-pe} or MIRA~\cite{Nolan2025-fu}), have enabled the identification of
diverse TCR pools with potential reactivity to specific antigens~\cite{Dash2017-if}, collected in
repositories like Immune Epitope Database (IEDB)~\cite{Vita2019-qn} and VDJdb~\cite{Shugay2018-nt}.
However, the current databases are unbalanced, with only a handful of dominant peptides (e.g.
derived from SARS-CoV-2, CMV, EBV, Influenza, etc.) and {MHC alleles} present, each linked to
hundreds of TCRs.

As a result, statistical inference and machine learning methods trained on such data to predict
TCR-pMHC reactivities lack generalizability beyond their training
data~\cite{Isacchini2021-fs,Weber2021-ik,Meynard-Piganeau2024-po}. On the structural side, only a
few hundred TCR-pMHC crystal structures have been experimentally generated to date. This is in part
because TCR-pMHC interactions are typically of low affinity and are difficult to stabilize to the
level needed for CryoEM or X-ray crystallography~\cite{Piepenbrink2009-jy}. Importantly, this
scarcity hinders computational prediction methods like AlphaFold in modeling reliable TCR-pMHC
structures~\cite{Bradley2023-rl,Yin2023-zm}. Machine learning models, such as large language models,
pre-trained on vast unannotated protein datasets can generate meaningful data
representations~\cite{Rives2021-sr,Lin2023-ub}. Fine-tuning these models on smaller, labeled
datasets for specific tasks (e.g., predicting mutation impacts on stability or function) has proven
highly effective~\cite{Madani2023-mg,Schmirler2024-wp,Wang2024-ep,Hie2024-vb}. Pre-trained
sequence-based protein language models have been employed to predict reactivities of antibodies to
antigens~\cite{Wang2024-ep,Hie2024-vb}, or TCRs to pMHC
complexes~\cite{Bravi2023-aj,Meynard-Piganeau2024-po}. However, lack of generalizability due to
sparse and skewed training data has remained a challenge in these models.

 
\begin{figure*}[th!]
\includegraphics[width=1.0\textwidth]{tcrpmhc/figures/schematic_0}
\figcaption{HERMES model for predicting and designing TCR-pMHC interactions.}{ {\bf (A)} HERMES is
a structure-based, rotationally equivariant neural network pre-trained on the protein universe to
predict amino acid preferences at a residue from its 3D structural
environment~\cite{Visani2024-sh,Pun2024-qz}. It takes as input all atoms within 10\AA\ of a focal
residue (with that residue’s atoms masked), featurizing them by atom type, charge, and SASA. HERMES
then outputs energies $E_i(\sigma)$ for the 20 possible amino acids $\sigma$ at residue $i$, which
serve as logits in a softmax to estimate amino acid likelihoods.
{\bf (B)} To predict TCR-pMHC affinity, we use two protocols: HERMES-{\em fixed} (left) and
HERMES-{\em relaxed} (right). In HERMES-{\em fixed} we use a fixed structural template closest to
the TCR-pMHC of interest as input to HERMES, and evaluate the peptide energy as the sum of
HERMES-computed per-residue energies for the peptide. In HERMES-{\em relaxed} mutations are
introduced to the template's peptide to match the peptide of interest, followed by local repacking
in PyRosetta~\cite{Chaudhury2010-mo}. Peptide energy is then evaluated similar to HERMES-{\em
fixed}, but using the relaxed structure. {\bf (C)} We also use two protocols for peptide design.
Similar to (B) we fix the closest structural template in HERMES-{\em fixed}, and use HERMES'
likelihood matrix for possible amino acids at each peptide position, which we sample to generate new
peptide sequences. In HERMES-{\em relaxed} (right), we begin with a random peptide $\sigma$, and
locally pack it with PyRosetta, and also generate an amino acid likelihood matrix (similar to
HERMES-{\em fixed}) for this packed structure. We sample a new sequence $\sigma'$ from this
likelihood matrix, and pack it with PyRosetta. We evaluate peptide energies
$E_{\text{peptide}}(\sigma)$, and $E_{\text{peptide}}(\sigma')$ for both peptides, using their
respective PyRosetta-relaxed structures, and we use these energies as the criterion to sample
improved sequences with MCMC (SI). This procedure is then iterated until convergence. Peptides
designed by both approaches are then filtered by TCRdock~\cite{Bradley2023-rl} or Alpha-Fold
3~\cite{Abramson2024-jp} to assure the stability of the complex.
\label{Fig:1}}
\end{figure*}
 
Protein-protein interactions in general, and immune-pathogen interaction in particular, rely on the
complementarity in the 3D structures and amino acid compositions of the interacting protein pairs.
Even though the interacting amino acids can be far apart in a protein sequence, they are close to
each other in structure, resulting in local interactions that can determine immune-pathogen
recognition. Machine learning models trained on protein structures instead of sequences can learn
local structural representations for amino acid statistics, which tend to be more generalizable
across different protein families and beyond their training
sets~\cite{Dauparas2022-fu,Blaabjerg2023-oq,Michalewicz2024-hw,Visani2024-sh,Pun2024-qz,Diaz2024-vv}.

Here, we present a structure-based approach to predict interaction affinities between TCRs and
peptides presented on {MHC class I}, and to design reactive peptides for specific TCR-MHC-I
complexes (TCR-MHC for short). Due to the limited availability of structural data for TCRs, we use
HERMES-- a physics-guided, 3D-equivariant machine learning model trained on the entire protein
universe to predict amino acid preferences based on their local structural
environments~\cite{Pun2024-qz,Visani2024-sh}.
 
{We demonstrate that HERMES accurately predicts both binding affinity of TCR-pMHC Class I complexes
and the functional responses and T-cell activities induced by a range of peptides presented on MHC
molecules. Affinity of TCR-pMHC complexes, measured as the dissociation constant $K_d$ by surface
plasmon resonance, is labor intensive to obtain. Functional readout used to measure T-cell
activities, including measurements of T-cell proliferation, labeled cellular activation markers, and
cytokine (e.g., IFN-$\gamma$) release, are easier to obtain than affinity measurements. Yet,
functional readouts are shaped not only by the TCR-pMHC affinity, but also by other cellular factors
such as the T-cell's differentiation state (e.g., naive or memory), self-reactivity of the
TCR~\cite{Altan-Bonnet2005-lm,Achar2022-tb}, and the avidity effects driven by TCR/co-receptor
clustering and pMHC density on antigen-presenting cells~\cite{Stone2009-ab}. Although HERMES is more
suited to predict molecular affinities, we show that its scores nonetheless can well track T-cell
activities, making it a reliable tool for predicting T-cell response to antigens.}

Lastly, we show that HERMES can be used to design novel peptides {binding specific} TCR-MHC
complexes, achieving up to 50\% experimental validation accuracy across diverse TCR-MHC systems.
%Our peptide design platform can facilitate the computational design of {neoantigen} libraries for
cancer vaccines.
By leveraging our design algorithm, we characterize the specificity of a diverse range of {TCRs},
providing a quantitative measure for the diversity of the peptide-MHC antigens that a TCR can
recognize in humans and mice.

\begin{figure*}[ht!]
	\centering
	\includegraphics[width=0.9\textwidth]{tcrpmhc/figures/affinity_figure_0}
	\figcaption{Predicting binding affinity of TCR-pMHC complexes.}{ {\bf(A)} The structure of the
	NY-ESO-1
	peptide (purple), in complex with HLA-A*02:01 and a purified specific 1G4 TCR is shown (PDB ID:
	2BNQ~\cite{Chen2005-jk}). This structure together with that of a peptide variant with C9V mutation
	(PDB ID: 2BNR~\cite{Chen2005-jk}) are used as templates to predict the binding affinity of the
	TCR-MHC complex to other peptide variants. {\bf (B)} The binding affinities ($-\log_{10}
	(K_d/\mu\text{M})$) of TCR-MHC complexes with 14 distinct peptides, measured using surface plasmon
	resonance, are shown against the predicted HERMES-{\em relaxed} peptide energies, using the closest
	structural template for each variant; 2BNQ is used for blue points and 2BNR for the orange point.
	The Spearman $\rho$ between predictions and measurements is reported in the panel. {\bf(C)} The
	heatmaps show the experimental affinity values (top) and predicted peptide energies (bottom) for
	all the peptide variants, with the identity of the mutations color coded based on the closest
	structural template. {\bf(D-F)} Similar to (A-C) but for the Tax viral peptide in complex with
	HLA-A*02:01 and a specific human TCR. Two structural templates are used~\cite{Garboczi1996-fr} for
	evaluating peptide energies with PDB IDs: 1AO7 (shown in D, and used for blue points), and 1QSF
	(used for orange points).
	\label{Fig:2}
	}
\end{figure*}

\section*{Model}

 T-cell response is mediated by the interactions between TCRs and pMHC complexes. To model the
 TCR-pMHC interaction, we use HERMES~\cite{Visani2024-sh}, a 3D rotationally equivariant
 structure-based neural network model, trained on the protein universe; see Fig.~\ref{Fig:1}A.
HERMES predicts the amino acid propensity from its surrounding 3D protein environment, a measure
that we have previously shown to serve as a reliable proxy for assessing the effect of mutations on
protein stability or binding~\cite{Pun2024-qz,Visani2024-sh}.

For TCR-pMHC complexes, we seek to determine how changes in a peptide sequence impact the binding
affinity of a TCR-pMHC complex, and ultimately, the T-cell response to the antigen. To do so, we
characterize a score, which we term peptide energy, for a peptide $\sigma$, given its surrounding
TCR and MHC structural environment,
{
\EQ
E_{\text{peptide}}(\sigma ; \text{TCR, MHC}) = \sum_{i=1}^\ell E(\sigma_i ; \text{TCR,
MHC},\sigma_{/i})
\label{eq:pepEn}
\EE
We assume that each peptide residue contributes linearly to the total energy by the amount
$E(\sigma_i ; \text{TCR, MHC}, \sigma_{/i})$, where $\sigma_i$ is the amino acid at position $i$ of
the peptide. A residue's energy contribution $E$ is evaluated by HERMES (logits of the network),
taking in as input the atomic composition of the surrounding TCR, MHC, and the rest of the peptide
$\sigma_{/i}$ (excluding the $i^{th}$ residue) within a 10~\AA~of $\sigma_i$'s $\alpha$-Carbon}; see
Fig.~\ref{Fig:1} and SI for details.

\begin{table*}[ht]
\centering
\begin{tabular}{ll|cc|cccccccc}
\toprule
 & & \multicolumn{2}{c|}{TCR-pMHC affinity} & \multicolumn{8}{c}{T-cell activity} \\
\cmidrule(lr){3-4}\cmidrule(lr){5-12}
Model type & Model & 1G4 TCR & A6 TCR & H2-scDb & TCR1 & TCR2 & TCR3 & TCR4 & TCR5 & TCR6 & TCR7 \\
\midrule
\multirow{1}{*}{Substitution Matrix} & BLOSUM62~\cite{Henikoff1992-jq} & \textbf{\underline{0.83}} &
{0.56}* & -0.05 & -0.07 & 0.14 & 0.38 & \textbf{\underline{0.33}} & 0.42* &
\textbf{\underline{0.55}} & -0.01 \\
\midrule
\multirow{2}{*}{Structure Prediction} & TCRdock~\cite{Bradley2023-rl} & {0.62*} &
\textbf{\underline{0.88}} & N/A & 0.02 & 0.13 & 0.21 & 0.18* & -0.02 & 0.19 & 0.06 \\
& TCRdock-{\em no template}~\cite{Bradley2023-rl} & 0.25 & -0.14 & N/A & -0.03 & 0.44 & 0.11 & -0.22
& 0.28* & 0.14 & 0.06 \\
\midrule
\multirow{1}{*}{Sequence based ML} & TAPIR~\cite{Fast2023-tk} & 0.29 & 0.35 & N/A & 0.36 & 0.36 &
-0.00 & 0.12 & 0.39* & 0.08 & -0.01 \\
\midrule
\multirow{2}{*}{Structure based ML} 
& ESM-IF1~\cite{Hsu2022-km} & 0.38* & 0.17 & 0.12 & 0.43* & 0.41 & 0.19 & -0.15 & 0.28* & 0.25 &
0.08 \\
& ProteinMPNN~\cite{Dauparas2022-fu} & -0.01 & 0.49 & \textbf{\underline{0.55}} & 0.46* & 0.52 &
0.49* & -0.28 & \textbf{\underline{0.45}} & 0.41* & \textbf{\underline{0.31}} \\
\midrule
\multirow{2}{*}{Structure based ML} & HERMES-\textit{fixed} & 0.29 & 0.63* & 0.46* &
\textbf{\underline{0.56}} & \textbf{\underline{0.71}} & \textbf{\underline{0.59}} & -0.06 & 0.37* &
0.42* & -0.01 \\
& HERMES-\textit{relaxed} & 0.72* & 0.56* & 0.29 & 0.31 & 0.57 & 0.55* & -0.20 & 0.37* & 0.42* &
0.14* \\
\bottomrule
\end{tabular}
\figcaption{Benchmarking of models for predicting TCR-pMHC binding affinities and T-cell
activities.}
{The table lists the Spearman~$\rho$ of models' predictions against experimentally
measured binding affinities (Fig.~\ref{Fig:2}) and T-cell activities (Fig.~\ref{Fig:3}) in response
to peptide variants. The HERMES-{\em fixed} and HERMES-{\em relaxed} performance are from models
with no noise (0.00) in training data, and the ProteinMPNN performance is from the model with a
small noise amplitude of 0.02~\AA; see SI for details. We show in \textbf{bold} the best model for
each dataset, and any model whose performance is \textit{not} significantly different from the best
model (p-value$>0.05$ associated with the difference in the Fisher's z-transformed Spearman
correlations) is indicated with an asterisk (*). See
Tables~\SpearmanrAffinityTable~and~\SpearmanrActivityTable~in the SI for $p$-values and additional
models.}
\label{table:spearmanr_affinity_and_activity}
\end{table*}


 To compute peptide energy, we input to HERMES an experimentally or computationally determined
 structure of a specified TCR-MHC complex bound to at least one peptide variant. Experimental data
 on the impact of peptide substitutions on the TCR-pMHC structure is limited, and computational
 models often fail to capture subtle conformational changes in the structure due to amino acid
 substitutions in a peptide. Given these constraints, we adopt two approaches to estimate peptide
 energies across diverse peptide variants for a given TCR-MHC complex:
\begin{itemize}
	\item HERMES-{\em fixed}: In the simplest approach, we choose a TCR-pMHC structure with the same
	TCR and MHC as our query and a peptide sequence closest to the peptide of interest. This structure
	is then used as our template in HERMES to compute the peptide energy as described in
	eq.~\ref{eq:pepEn} (Fig.~\ref{Fig:1}B). This method does not alter the underlying structure and
	assumes that peptide amino acid substitutions do not significantly change the conformation of the
	TCR-pMHC complex.
	
	\item HERMES-{\em relaxed}: Since amino acid substitutions can locally reshape a protein complex,
	we introduce a more involved protocol to account for these structural changes. We begin with the
	closest available structure for the TCR-pMHC complex. We then mutate the original peptide to the
	desired state and apply a relaxation procedure using PyRosetta to pack the substituted amino
	acid~\cite{Chaudhury2010-mo} (SI). During relaxation, side-chain and backbone atoms of the peptide
	are allowed to move, while only the side-chain atoms of the MHC and TCR chains within a
	10~\AA\,radius of the peptide are flexible. We then calculate the peptide energy with HERMES
	(eq.\ref{eq:pepEn}), using the relaxed structure as input (Fig.~\ref{Fig:1}B). Since PyRosetta's
	relaxation procedure is stochastic, we average the peptide energy across 100 realizations of these
	relaxations. {For ablation purposes, we also include results obtained by selecting the relaxed
	conformation with the lowest (most favorable) Rosetta energy, and using its HERMES peptide energy.}
\end{itemize}





\section*{Results}
\subsection*{Predicting binding affinity of TCR-pMHC complexes} 
 Binding between TCRs and pMHC complexes is necessary to mediate a T-cell response. The binding
 affinity of a natural TCR-pMHC complex is relatively low, with a dissociation constant $K_d$ in the
 range of $1-100\mu{\rm M}$~\cite{Davis1998-bv,Piepenbrink2009-jy,Zhong2013-ja}, in contrast to
 nanomolar affinities for antibody-antigen complexes. Engineered TCRs can achieve much higher
 affinities in a nanomolar~\cite{Soto2013-df} to picomolar~\cite{Li2005-md, Dunn2006-dh} range.

Surface plasmon resonance (SPR) spectroscopy provides accurate measurements of the dissociation
constant ($K_d$) for specific TCR-pMHC systems. In these experiments, one protein--either the TCR or
the MHC--is immobilized on a conducting plate, and the other is introduced in solution to bind to
it. This binding alters the local refractive index near the plate, affecting the resonance signal.
We {used a published} dataset of SPR-measured affinities for two TCR-MHC complexes binding to an
ensemble of peptides, for which crystal structures of the complexes with at least one peptide are
available~\cite{Aleksic2010-wz,Pettmann2021-kl}; see Table~\AffActInfoTable~for details.

In our first example, we examined TCR 1G4 binding to NY-ESO-1 peptide variants, which is a
cancer-testis antigen commonly expressed in many cancers and is targeted by
immunotherapies~\cite{Esfandiary2015-jf}. We used two structural templates, one with the wild-type
peptide SLLMWITQC (NY-ESO-1), and the other with a more immunogenic peptide, in which the Cysteine
at position 9 is substituted by Valine (C9V)~\cite{Chen2000-vq,Chen2005-jk} (Fig.~\ref{Fig:2}A). For
variants differing by a single amino acid, our predictions with HERMES-{\em relaxed} correlated well
(Spearman {$\rho=0.72$}) with experimentally measured affinities from
refs.~\cite{Aleksic2010-wz,Pettmann2021-kl}; see Fig.~\ref{Fig:2}B,~C.

Next, we tested the A6 TCR binding to the HTLV-1 Tax peptide LLFGYPVYV. We employed two structural
templates: one with the wild-type, the other with the mutant peptide with substitution
Y8A~\cite{Garboczi1996-fr} (Fig.~\ref{Fig:2}D and {Table~\AffActInfoTable}). Again, the HERMES-{\em
relaxed} peptide energies correlated well (Spearman {$\rho=0.56$}) with the affinities, measured in
ref.~\cite{Pettmann2021-kl} (Fig.~\ref{Fig:2}E,~F).

{Overall, HERMES predictions aligned closely with these limited experimental affinity measurements,
with HERMES-{\em relaxed} showing a better agreement with affinity measurements than HERMES-{\em
fixed}.
We compared our approach for scoring peptides to alternative methods, including the baseline
substitution matrices (BLOSUM62~\cite{Henikoff1992-jq} and the TCR-specific substitution matrices
from ref.~\cite{Luksza2022-gx}), the score from the TCR-pMHC structure prediction algorithm
TCRdock~\cite{Bradley2023-rl}, the sequence-based machine learning algorithms for TCRs
(TAPIR~\cite{Fast2023-tk}, NetTCR~\cite{Jensen2024-xs}, and TULIP~\cite{Meynard-Piganeau2024-po}),
and the structure-based inverse folding models (ESM-IF1~\cite{Hsu2022-km} and
ProteinMPNN~\cite{Dauparas2022-fu}); see Table~\ref{table:spearmanr_affinity_and_activity} for key
benchmarking results and Table~\SpearmanrAffinityTable~and Fig.~\ProteinMPNNScoringSchemesFigure~for
a more extensive comparison of different methods, including the different HERMES models, and
different scoring schemes by ProteinMPNN.
 
For the 1G4 TCR, BLOSUM62~\cite{Henikoff1992-jq} achieves the highest correlation with experimental
affinity data, whereas template-guided TCRdock~\cite{Bradley2023-rl} leads for the A6 system; in
contrast, TCRdock without using a homologous structure as a template underperforms. In both systems,
HERMES achieves performance that is statistically indistinguishable from the top performers
(Fisher's z-test p-value $>0.05$). Overall, BLOSUM62, template-guided TCRdock, and HERMES rank
strongly as predictors of TCR-pMHC affinities across the tested datasets. In addition, for the 1G4
system, ESM-IF1~\cite{Hsu2022-km} and the TCR-specific substitution matrices from
ref.~\cite{Luksza2022-gx} also perform competitively. However, due to the limited number of affinity
measurements, we lack statistical power to resolve finer differences among the models.

It should be noted that although BLOSUM62 is one of the best contenders for affinity prediction, its
sequence‑averaged construction ignores local structural context and therefore assigns identical
substitution scores regardless of a residue's environment (Fig.~\HsiueHisFigure); similar
limitations apply to the TCR-specific substitution matrices from ref.~\cite{Luksza2022-gx}. HERMES,
by contrast, is structure‑aware and can attribute distinct scores to the same mutation when it
occurs in different conformational settings--an ability that is critical for nuanced mutational
forecasts (Fig.~\HsiueHisFigure). Determining the biophysical circumstances under which such
context‑agnostic scoring schemes suffice, and when structure‑aware approaches become necessary, will
require systematic evaluation on broader, more diverse datasets.


Lastly, we probed HERMES's robustness to potential inaccuracies in structures generated by the
widely used AlphaFold3 (AF3) algorithm~\cite{Abramson2024-jp}, replacing the crystal structure
inputs with AF3 models; see SI for details.
Both the template-guided and the de novo AF3 structures remained within $1.1$~\AA\, RMSD of the
respective crystal structures (Table~\AfMetricsTable).
For both HERMES and ProteinMPNN, replacing experimental structures with AF3 predictions altered
predictions for the 1G4 and A6 TCRs binding to pMHC variants, but the changes were not statistically
significant (Fisher's z-test p-value $>0.05$); see Fig.~\AfScoringResultsFigure. It is nonetheless
expected that the predictions should suffer in the absence of reliable structural templates, as seen
also in other methods such as TCRdock~\cite{Bradley2023-rl}
(Tables~\ref{table:spearmanr_affinity_and_activity},~\SpearmanrAffinityTable). }
 
 
\subsection*{\bf Predicting T-cell response to peptide antigens}
{T-cell activation is shaped by factors beyond TCR-pMHC affinities, including the T-cell's
differentiation state, its cross-reactivity to
self-antigens~\cite{Altan-Bonnet2005-lm,Achar2022-tb}, and avidity effects from TCR/co-receptor
cross-linking~\cite{Stone2009-ab}. Nonetheless, because readouts of T-cell activiy (e.g., T-cell
proliferation, and production of activation markers) are easier to obtain than SPR-based affinity
measurements, we next asked whether HERMES scores, which are more attuned for predicting TCR-pMHC
binding affinity, can serve as predictors of T-cell activities.}

As the first system, we examined T-cell responses to all single-point mutants of a native peptide,
as measured in ref.~\cite{Luksza2022-gx} for different TCR-pMHC systems. Specifically, we make
predictions for (i) three TCRs (TCRs~1-3) recognizing variants of a highly immunogenic
HLA-A*02:01-restricted human cytomegalovirus (CMV) epitope NLVPMVATV (NLV), (ii) three TCRs
(TCRs~4-6) recognizing variants of the {weaker} HLA-A02:01-restricted melanoma gp100 self-antigen
IMDQVPFSV, and (iii) a TCR (TCR~7) recognizing variants of weakly immunogenic HLA-B*27:05-restricted
pancreatic cancer neo-peptide GRLKALCQR, believed to have elicited an effective immune response in
long-term cancer survivors~\cite{Luksza2022-gx}. T-cell activity was quantified by determining the
half-maximal effective concentration (EC$_{50}$) from the dose-response curves of 4-1BB$^+$ CD8$^+$
T-cells across varying peptide concentrations; see Methods on details of inferring EC$_{50}$.

With the exception of TCR1 (where a TCR-pMHC template with HLA-A*02:01 and one peptide variant was
already available), we relied on AlphaFold3 (AF3)~\cite{Abramson2024-jp} to model the remaining
TCR-pMHC complexes; see Table~\AffActInfoTable~and Dataset~\AfInfoDataset~for template information
and AF3 sequence inputs for TCRs~2-7. {As shown in Figs.~\ref{Fig:3}A,
and~\ActivityHeatmapsFigure,~HERMES-{\em fixed} peptide energy predictions (eq.~\ref{eq:pepEn})
correlate well with experimental EC$_{50}$ measurements for the CMV peptide variants, reaching up to
Spearman correlation $\rho= 0.71$ for TCR2, and outperforming the alternative models. For TCR1 and
TCR3, the performance of ProteinMPNN~\cite{Dauparas2022-fu} is worse but statistically
indistinguishable from HERMES
(Tables~\ref{table:spearmanr_affinity_and_activity},~\SpearmanrActivityTable, and
Fig.~\ProteinMPNNScoringSchemesFigure). Moreover, replacing the crystal structure for TCR1 with an
AF3 model yields comparable performance, underscoring HERMES's robustness to structural inputs so
long as the structural models are of high quality
(Figs.~\AfExamplesFigure,~\AfScoringResultsFigure). Correlations are lower for TCRs~4-6 targeting
variants of the melanoma gp100 self-antigen, with HERMES achieving Spearman correlation
{$\rho=0.42$} for TCR6 as its best prediction in this class, with an accuracy comparable to
ProteinMPNN~\cite{Dauparas2022-fu} (Fig.~\ref{Fig:3}B {and
Tables~\ref{table:spearmanr_affinity_and_activity},~\SpearmanrActivityTable}).}
HERMES underperforms for TCR7 responding to the pancreatic cancer neo-peptide GRLKALCQR
(Fig.~\ref{Fig:3}C){, for which only the ProteinMPNN's predictions show significant albeit weak
correlation with data (Table~\ref{table:spearmanr_affinity_and_activity})}. The reduced accuracy of
HERMES for cancer epitopes compared to the variants of the viral CMV epitope could reflect the
impact of antagonism from self-antigens similar to tumor epitopes, hindering T-cell
responses~\cite{Altan-Bonnet2005-lm,Achar2022-tb}---though additional experiments are needed to
confirm this.

{Sequence-based models (TAPIR~\cite{Fast2023-tk}, NetTCR~\cite{Jensen2024-xs}, and
TULIP~\cite{Meynard-Piganeau2024-po}) are outperformed by structure-based models tested here in all
cases except for TCR4, which has the lowest AF3 confidence at the TCR-pMHC interface (i.e., a lower
confidence in its fold); see Table~\AfMetricsTable~and Fig.~\AfScoringResultsFigure~for the AF3
confidence metrics for all systems. Here, even the basic substitution models
(BLOSUM62~\cite{Henikoff1992-jq} and the TCR-specific matrices from ref.~\cite{Luksza2022-gx})
outperform the structure-based predictors
(Tables~\ref{table:spearmanr_affinity_and_activity},~\SpearmanrActivityTable). This suggests
that structural approaches hinge on a reliable structure as input, though more data is needed to
draw definitive conclusions (Fig.~\AfScoringResultsFigure).
}

Next, we examined the {therapeutic T-cell responses} to all single-point mutants of the
HLA-A*02:01-restricted p53-derived {neoepitope} p53$^{\rm R175H}$, HMTEVVRHC, measured in
ref.~\cite{Hsiue2021-ag}. This T-cell therapy approach used a designed bi-specific Fab antibody that
binds both to the TCR-CD3 complex on T-cells and to the variants of the p53$^{\rm R175H}$-MHC
complex on tumor cells (Table~\AffActInfoTable). A strong binding to an {epitope} activates the
engaged T-cell, and the extent of such activity was measured by the production of IFN-$\gamma$. The
resulting HERMES-{\em fixed} predicted peptide energies (eq.~\ref{eq:pepEn}) show a {Spearman
$\rho=0.46$} with the experimentally measured IFN-$\gamma$ production
(Figs.~\ref{Fig:3}D,~\ActivityHeatmapsFigure). {ProteinMPNN achieves a slightly higher correlation
(Spearman $\rho=0.55$), but this improvement over HERMES is not statistically significant (Fisher's
z-test p-value $>0.05$); see Table~\ref{table:spearmanr_affinity_and_activity}. Lastly, we found
that replacing the crystal structure with the AF3 model substantially degrades HERMES and
ProteinMPNN performances (Fig.~\AfScoringResultsFigure,
Tables~\ref{table:spearmanr_affinity_and_activity},~\SpearmanrActivityTable), as the AF3 model
consistently mis-docks the bispecific antibody on the pMHC complex (Table~\AfMetricsTable~and
Fig.~\AfExamplesFigure).
}
 


\begin{figure*}[ht!]
	\centering
	\includegraphics[width=0.9\textwidth]{tcrpmhc/figures/activity_all_0}
	\figcaption{Predicting T-cell activity in response to peptide variants.}{
	{\bf (A-C)} The predicted HERMES-{\em fixed} peptide energies versus experimentally measured T-cell
	activity are shown for (A) HLA-A*02:01-restricted CMV peptide variants interacting with TCR1-TCR3
	(panels), (B) HLA-A*02:01-restricted melanoma gp100 peptide variants interacting with TCR4-TCR6,
	and (C) HLA-B*27:05-restricted pancreatic cancer peptide variants interacting with TCR7. In each
	case, T-cell activity in response to all single amino acid mutations from the wild-type is
	measured. T-cell activity is based on fitted EC$_{50}$ values from dose-response curves of
	4-1BB$^+$ CD8$^+$ T-cells (SI; data from \cite{Luksza2022-gx}). Gray points indicate peptides with
	T-cell responses below detection limit (no EC$_{50}$ obtained) and are {excluded from the reported
	Spearman correlations.} Except for TCR1 (resolved structure, PDB ID: 5D2N~\cite{Yang2015-cj}), the
	other TCR-pMHC complexes were modeled with AF3, and the model's peptide PAE for each wild-type
	peptide is denoted. These structures are used as templates to estimate peptide energies with
	HERMES-{\em fixed}; see Table~\AffActInfoTable~and Dataset~\AfInfoDataset~for details on the
	templates. In all panels, TCR-$\alpha$ (TRA) is shown in blue, TCR-$\beta$ (TRB) in orange, the
	peptide in magenta, and HLA in gray. {\bf (D)} The predicted HERMES-{\em fixed} peptide energies
	are shown against experimentally measured T-cell activity for single point mutants of the
	HLA-A*02:01-restricted p53$^{\rm R175H}$ {neoepitope} from ref.~\cite{Hsiue2021-ag}. T-cell
	response (IFN-$\gamma$ production) is mediated through a bispecific antibody binding to the pMHC
	complex. The co-crystallized structure of the bispecific antibody with the p53$^{\rm
	R175H}$-HLA-A*02:01 complex (PDB ID: 6W51~\cite{Hsiue2021-ag}) is used as the template for
	HERMES-{\em fixed} predictions. For all systems (A-D), the heatmaps for the predicted and the
	experimentally measured mutational effects are shown in Fig.~\ActivityHeatmapsFigure.
		\label{Fig:3}
	}
\end{figure*}



Across the limited number of systems examined, our analyses indicate that--when a suitable
structural template is available--HERMES can estimate epitope mutational effects at the pMHC-TCR
interface and, in the H2-scDb case, at a Fab-pMHC interface. In this benchmark, ProteinMPNN performs
comparably to HERMES, indicating that inverse-folding algorithms not explicitly trained for this
task may nonetheless capture the impact of peptide mutations on T-cell responses. However,
definitive conclusions about the generalizability of these methods to neoepitopes will require
larger datasets with broader mutational coverage across peptide variants. Importantly, predictive
performance for both approaches depends on access to a high-quality structure (experimental or
computational), which limits applicability in systems without suitable structural models
(Fig.~\AfScoringResultsFigure); 
see SI and Tables~\ref{table:spearmanr_affinity_and_activity}~and~\SpearmanrActivityTable~for a
detailed benchmark of different approaches, including different HERMES models.


%Our analysis shows that, despite not being explicitly trained for this purpose, HERMES peptide
energy effectively predicts T-cell responses to {neoepitopes} presented by different MHC-I alleles.
Moreover, HERMES predictions are not limited to TCR-pMHC interactions but also extend to
antibody-pMHC interactions, commonly leveraged in T-cell therapies for their higher achievable
affinities;


\begin{figure*}[th!]
\centering
\includegraphics[width=0.63\textwidth]{tcrpmhc/figures/design_figure_final_4_0}
\figcaption{{De novo} peptide design to elicit a T-cell response.}{ {\bf (A-C)} Validation results
are
shown for designed peptides built upon the wild-type templates from (A) NY-ESO peptide against TCR
1G4 and HLA-A*02:01 (PDB IDs: 2BNR, 2BNQ~\cite{Chen2005-jk}), (B) EBV epitopes against the TK3 TCR
and the HLA-B*35:01 (PDB IDs: 3MV7~\cite{Gras2010-xp}, 4PRP~\cite{Liu2014-cz}), and (C) MAGE-A3 and
Titin epitopes against an engineered TCR and HLA-A*01:01 (PDB IDs: 5BRZ, 5BS0~\cite{Raman2016-ew}).
Designs were validated using Jurkat cells {with endogenous NFAT-eGFP reporter} expressing the paired
TCR, interacting with peptide presented by the MHC in each system. The reported T-cell activities
(percentage GFP-positive) are averaged over three replicates. The box plots show T-cell activity for
designed peptides at different Hamming distances from the wild-type; the black line is the median,
the box spans the 25-75\% quantiles, and dots indicate individual data points. Red and pink
boxes/markers correspond to designs with HERMES-{\em fixed} and HERMES-{\em relaxed} models,
respectively; square markers denote wild-type peptides, and stars are negative controls (no
peptide). The dashed line marks the activation threshold, and numbers below each bar show the
relative number of successful designs out of total designs. The top plot in each panel shows the
size of the peptide sequence space at different distances from the wild-types, with the gray line
indicating the complete sequence space, and the green indicating the typical number of sequences
likely to be presented by the MHC of each system; see SI for details. The bottom shows sequence
motifs as PWMs for (i) true-positive designs that induced T-cell activation, (ii) false positives
that did not, and (iii) the MHC presentation motif from the MHC Motif Atlas~\cite{Tadros2023-ij}.
{\bf (D)} 3D structures of TCR–pMHC complexes focused on the peptide are shown for MAGE-A3 (PDB ID:
5BRZ), Titin (PDB ID: 5BS0), and the highest-activity design in (C) modeled by
AF3~\cite{Abramson2024-jp}. The enhanced activity for Titin versus MAGE-A3, and for the designed
peptide versus both wild-types follows from the formation of polar bonds between position 8 of the
peptide (magenta) and the CDR3 loop of TCR$\beta$ (orange). The PWM highlights the prominence of
E$_8$ among the eight designs with improved T-cell activity in (C). {{\bf (E)} Structures predicted
with AF3 for the activity-inducing designed peptides (above the dashed lines in A-C) are displayed
in complex with their cognate TCR-MHC's. For each system, designs generated by {\em HERMES-fixed}
(red) and {\em HERMES-relaxed} (pink) are overlaid after aligning the MHC molecule (in gray). Slight
shifts are seen in the backbone conformations of the peptides and the TRA (blue) and TRB (orange)
loops.}
\label{Fig:4}}
\end{figure*}


\subsection*{\bf Design of novel peptides reactive to a TCR-MHC complex} 
TCR-engineered T-cell therapies are designed to target defined neoantigens, but specificity
remains a central challenge. Engineered high-affinity TCRs can cross-react with untested
self-peptides, causing severe off-target toxicity; in a MAGE-A3 directed trial, recognition of a
Titin-derived peptide led to fatal cardiac events~\cite{Linette2013-tc}. Therefore, developing
approaches to predict the distribution of peptides recognized by engineered TCRs could enable more
stringent preclinical screening and help mitigate off-target risk in patients.


To map the peptide-recognition landscape of TCRs, we present a structure-guided pipeline for
designing immunogenic peptides predicted to elicit strong TCR recognition when presented by MHC
class I. Leveraging HERMES's ability to predict T-cell responses to peptides presented by MHC-I, our
approach begins with an existing template structure--experimentally or computationally resolved---of
a TCR-MHC in complex with at least one peptide variant, which we refer to as the wild-type peptide.
We then generate candidate peptides using two strategies, with different degrees of complexity; see
Fig.~\ref{Fig:1}C for a schematic description of the two pipelines.

In the basic pipeline (HERMES-{\em fixed}), we sequentially mask the atoms of each amino acid along
the peptide, one at a time.
Using the HERMES model, we predict the probability of each amino acid type for a given residue,
based on the local structural environment within 10~\AA\, of the masked residue. Repeating this
procedure for every amino acid along the peptide yields a position weight matrix (PWM) that
represents the probabilities of different amino acids at each peptide position. Peptides are then
sampled by drawing from this PWM. It should be noted that HERMES-{\em fixed} retains the structural
fingerprints of the original peptide, as it does not account for local structural changes induced by
amino acid substitutions; see Fig.~\ref{Fig:1}C and~\HermesDesignsEnergiesFigure A, and SI for more
details.

To incorporate structural flexibility and relaxation following amino acid substitutions, we
introduce the design pipeline of HERMES-{\em relaxed}, using simulated annealing and Markov chain
Monte Carlo (MCMC) sampling. We begin with the template TCR-MHC structure and a completely random
peptide sequence, which is packed {\em in silico} within the structure using PyRosetta's packing
functionality, followed by the Relax protocol~\cite{Chaudhury2010-mo}. During relaxation, both
side-chain and backbone atoms of the peptide are allowed to move, while only the side-chain atoms of
the MHC and TCR chains within 10~\AA~of the peptide are flexible. With this relaxed structure, we
apply the HERMES-{\em fixed} method to generate a PWM and sample a new peptide, which is then packed
and relaxed to form a new structure. We use MCMC to determine whether to accept the newly sampled
peptide by comparing its HERMES peptide energy (favorability), evaluated in its relaxed
conformation, to that of the original peptide (eq.~\ref{eq:pepEn}). We then iteratively perform this
procedure, incrementally reducing the MCMC temperature at each step, until the results converge; see
Fig.s~\ref{Fig:1}C,~\HermesDesignsEnergiesFigure B-C,~and SI for more details.

The peptides designed by both approaches are then filtered by TCRdock~\cite{Bradley2023-rl} or
{AlphaFold3}~\cite{Abramson2024-jp}, using the Predicted Alignment Error (PAE) between TCR and pMHC
interfaces, to assure folding and binding of the resulting TCR-pMHC structures. For each system, we
define a specific PAE threshold close to the TCRDock PAE of the system's TCR-MHC in complex with the
native (wild-type) peptide, which is known a priori to activate T-cells. As a result, the PAE
thresholds vary slightly across different systems; see SI for details.

We tested the accuracy of our peptide design pipeline in three systems, with the structural
templates taken from (i) the cancer-testis antigen NY-ESO-1 in complex with the TCR 1G4 and
HLA-A*02:01, (ii) a peptide derived from the Epstein-Barr virus (EBV) in complex with the TK3 TCR
and the HLA-B*35:01, and (iii) the immuno-therapeutic target MAGE-A3 in complex with an engineered
TCR and HLA-A*01:01. For brevity, we refer to these three systems as NY-ESO, EBV and MAGE,
respectively; See Table~\DesignInfoTable~for information on the structure templates used in each
case. {These systems were selected because each one has high-resolution crystal structures of the
same TCR-MHC scaffold bound to at least two distinct peptide variants, and that they span a
therapeutically relevant spectrum that includes both engineered (for MAGE and NY-ESO) and natural
(for EBV) TCRs directed against viral as well as tumor antigens.} It should be noted that as we
progressed from NY-ESO to EBV and MAGE, we chose to explore the sequence space farther from the
wild-type peptides, leading to more challenging designs.

We validated our designs using Jurkat cells expressing {a specific} TCR {and an endogenous}
NFAT-eGFP reporter to indicate T-cell activation in response to different peptides. We tested the
TCR-MHC specificity in all systems--NY-ESO, MAGE, EBV--under 96 different conditions, including the
{de novo} designed peptides, positive controls (wild-type peptides from the template structures),
and an unstimulated ``no peptide" control. Peptides were presented on artificial APCs (aAPCs),
expressing the specified HLA in each system. For NY-ESO designs, we determined the {the percentage
of GFP-positive cells} through both flow cytometry and fluorescence microscopy. Given the
consistency between the two experimental approaches (Fig.~\FlowMicroscopyFigure) we relied on
fluorescence microscopy only to measure GFP levels induced by different peptides for MAGE and EBV
designs; see SI for experimental details and Dataset~\experimentalMeasurementsDataset~for data on
these measurements.

We used two structural templates for the NY-ESO system: one containing the original 9-amino acid
peptide SLLMWITQC, and another with the more immunogenic C9V
substitution~\cite{Chen2000-vq,Chen2005-jk}. Our {de novo} designs differed from the closest
wild-type peptide by 2 to 7 amino acids, with a TCRDock PAE~$\leq 5.5$. For this system only, we
selected 35 negative designs (designed by HERMES, but with PAE $> 5.5$) for experimental testing:
only one of these was a false negative, which interestingly, had a desirable AF3 PAE score (as low
as the wild-type peptides); see Figs.~\DesignTCRDockAFScatterFigure,~\DesignExpNyesoFigure. For the
58 positive designed peptides (PAE~$\leq~5.5$), we achieved an overall design accuracy of 50\%, with
GFP levels significantly higher than the negative control; see SI for details. The accuracy of
predictions decreases as the sequence divergence of the {de novo} designs from the wild-type
peptides increases (Fig.~\ref{Fig:4}A,~\DesignExpNyesoFigure). HERMES-{\em fixed} achieved higher
design accuracy of 70\% within its smaller designed sequence subspace of up to 5 substitutions from
the wild-types, whereas HERMES-{\em relaxed} explored sequences further from the wild-type peptides
(4-7 substitutions) but with a reduced accuracy of 24\% (Fig.~\ref{Fig:4}A). The main difference
between the true and false positive sequences was at position 8--from strongly preferring glutamine
within the true positives, to having similar preferences among glutamine, glycine, arginine, or
methionine within the false positives (Fig.~\ref{Fig:4}A).

For the EBV system, we used two structural templates: one with the 11-amino acid peptide
HPVGEADYFEY---commonly referred to as HPVG---derived from the viral latent nuclear protein EBNA-1,
and the other with an epitope from a wide-spread viral {variant}, with a glutamine at position
5~\cite{Liu2014-cz}. Unlike NY-ESO, the HPVG peptide has a more flexible conformation with a helical
turn within the {MHC} binding pocket. Our {de novo} designs differed from the closest wild-type
peptide by 3 to 9 amino acids, with a TCRDock PAE~$\leq~5.1$. Across the 93 {de novo} designs, we
achieved an overall accuracy of 14\%, associated with the peptides that induced significant GFP
levels in T-cells. The accuracy of predictions decreases with increasing the sequence divergence
from the wild-types (Fig.~\ref{Fig:4}B,~\DesignExpEbvFigure). Notably, we achieved 40\% accuracy
among the designs within 3-5 amino acid distance of one of the wild-type templates, while none of
the designs with larger than 5 amino acid distance were successful. Lastly, we observe a strong
preference for {g}lutamic acid at position 10 of the successful designs relative to the false
positives (Fig.~\ref{Fig:4}B).


\begin{figure*}[th!]
\includegraphics[width=0.9\textwidth]{tcrpmhc/figures/entropy_TCRMHC_cartoon_0}
\figcaption{Evaluating T-cell specificity from {de novo} designed peptides.}{ {\bf (A)} 
T-cell specificity is measured by peptide entropy, which is computed from the HERMES-{\em fixed}
predicted amino acid likelihood matrix, starting from a TCR-pMHC structural template. Each point on
the plot corresponds to the entropy estimate from a distinct TCR-pMHC template, grouped by their MHC
allele. The box plots display the distribution of peptide entropies for TCR-pMHC templates sharing
the same MHC allele, with a white line marking the median and the box spanning the 25\% to 75\%
quantiles. The number on the side of each box indicates the number of distinct structural templates
we used for that MHC class; see Table~\AllPdbsTable~for details on these structures. The side plot
shows peptide degeneracy for MHC presentation alone, evaluated as the entropy of the peptide
position weight matrices gathered from the MHC Motif Atlas~\cite{Tadros2023-ij}. The error bars
indicate the standard deviation of peptide degeneracies across motifs of different lengths,
associated with a given MHC allele. {\bf (B)} Similar to (A) but for degeneracy of peptides
interacting with TCRs and MHCs in mice. {\bf (C)} Polar bonds (blue dashed lines) between the NY-ESO
peptide (magenta) and HLA-A*02:01 (gray) is shown at the top (PDB ID: 1S9W~\cite{Webb2004-am}). The
additional polar bonds formed between the peptide and TRA (blue) and TRB (orange) of an interacting
TCR are shown at the bottom (PDB ID: 2BNR~\cite{Chen2005-jk}). These additional polar bonds limit
the degeneracy of the peptides that can interact with a TCR-MHC complex compared to those presented
by the same MHC.
\label{Fig:5}
}
\end{figure*}

For the MAGE system, we used two structural templates: one with the 9-amino acid MAGE-A3 peptide
EVDPIGHLY and another with the Titin-derived self-peptide ESDPIVAQY, which is expressed in
cardiomyocytes~\cite{Raman2016-ew}. %The Titin-derived peptide is known to have triggered an
off-target response of the affinity-enhanced MAGE-A3-specific engineered TCR that resulted in severe
cardiac toxicity and deaths of two patients~\cite{Linette2013-tc}.
Our {de novo} designs differed from the closest wild-type peptide by 3 to 8 amino acids, with a
TCRDock PAE~$\leq~5.2$ (Fig.~\ref{Fig:4}C). From the resulting 93 {de novo} designs, 20\% activated
T-cells with significant levels of GFP expression. Similar to the other systems, the accuracy of
predictions decreases with increasing the sequence divergence from the wild-types
(Fig.~\ref{Fig:4}C). We achieved 30\% accuracy among the designs with 3-5 to amino acid distance,
while none of the designs with larger than 5 amino acid distance from the wild-types were
successful.

Fig.~\ref{Fig:4}C shows differences in amino acid composition between successful and failed designs.
Alanine and glycine scanning experiments have previously suggested that MAGE-A3 peptide variants
with an E-DPI-\,-\,-Y motif activate T-cells in this system~\cite{Cameron2013-vw}. Consistent with
these findings, the fraction of T-cell activating variants in our designs increases as their
sequences more closely resemble this motif (Fig.~\DesignExpMageFigure).

In a subset of designs using the MAGE-A3 template, we fixed the glutamic acid at position 1 (E$_1$),
which is part of the E-DPI---Y pattern. In both MAGE-A3 and Titin E$_1$ is within 4~\AA~of the MHC
and forms strong polar interactions with it~(Fig. 4D). Fixing E$_1$ improved the design accuracy to
$1/3$, compared to $1/9$ when this amino acid was not constrained. Even without this constraint in
the design protocol, E$_1$ appeared in $2/3$ of successful designs, and in none of the unsuccessful
ones (Fig.~\DesignExpMageFigure). Notably, among the more distant designs made with HERMES-{\em
relaxed}, only those containing a fixed E$_1$ were able to activate T-cells (Fig.~\ref{Fig:4}C).
This suggests that constraining essential amino acids can enable broader exploration of sequence
space at other positions, possibly due to epistatic interactions. However, further investigation is
needed to confirm this.

{Several} of our designs outperformed the wild-type peptides in activating T-cells: eight surpassed
Titin’s GFP levels, and eleven exceeded MAGE-A3’s (Fig.~\ref{Fig:4}C). These improved designs
strongly favor glutamic acid (E) at position 8 of the peptide, where MAGE-A3 has leucine (L), and
Titin has glutamine (Q) (Fig.~\ref{Fig:4}D). In Titin, the glutamine forms a polar bond with an
arginine (R) in the interacting TCR$\beta$, and it is likely responsible for increasing Titin's
activity relative to MAGE-A3's. The glutamic acid substitution in our designs further strengthens
this polar bond (Fig.~\ref{Fig:4}D), resulting in a stronger T-cell response. {For NY-ESO and EBV,
the wild-type peptides were strongly immunogenic and induced high levels of GFP, and many of our
designs induced comparable T-cell responses, but did not improve upon those (Fig.~\ref{Fig:4})}.


{To gauge the impact of using AF3 structures as design templates, we reran the HERMES-{\em fixed}
protocol with AF3 predicted structures (with and without template guidance), and obtained peptide
PWMs that were almost identical to those derived from using crystal structures (SI and
Fig.~\DesignPWMsWithAf). An interesting exception is the MAGE system, where using the AF3 template
introduced a strong preference for glutamic acid at positing 1 of the peptide
(Fig.~\DesignPWMsWithAf C), mirroring the pattern observed among the designed peptides that induced
strong T-cell response (Fig.~\ref{Fig:4}D).}

Although no direct experimental comparisons were performed, computational evaluations using TCRdock
PAE and presentation score by MHC-I alleles indicate that HERMES designs outperform or match those
generated by ProteinMPNN~\cite{Dauparas2022-fu}, the corresponding MHC position weight
matrix~\cite{Lundegaard2008-ao, Lundegaard2008-wz, Andreatta2016-uh,Tadros2023-ij}, or the BLOSUM
substitution matrix~\cite{Henikoff1992-jq} (SI and
Figs.~\DesignTcrdockNyesoFigure,~\DesignTcrdockEbvFigure,~\DesignTcrdockMageFigure). A notable
exception is in the MAGE system, where ProteinMPNN designs generally have more favorable PAE scores,
albeit with a lower proportion of designs passing the antigen presentation filter by the MHC
molecule (Fig.~\DesignTcrdockMageFigure).

In all the three systems we used the PAE values from TCRdock and AF3 as filtering criteria to select
designs for experimental validation (see SI). The PAE values reported by the structure prediction
algorithms have been shown to be noisy indicators of protein folding fidelity and design
quality~\cite{Roney2022-am}. While larger experimental libraries are needed to draw definitive
conclusions, it appears that PAEs from TCRdock and AF3 have some complementary predictive strengths
({Figs.~\DesignTCRDockAFScatterFigure,~\DesignAUCFigure}), and therefore, combining their
information could lead to a more robust design decisions (SI). {Both of these scores correlate only
weakly
with HERMES
(Figs.~\DesignTCRDockAFScatterFigure,~\DesignTcrdockNyesoFigure,~\DesignTcrdockEbvFigure,~\DesignTcrdockMageFigure),
so they can provide complementary design filters. Furthermore, designed peptides with higher (more
favorable) HERMES energies more often activated T-cells (Fig.~\DesignAUCFigure), suggesting that a
more restrictive sampling (e.g., with a lower MCMC temperature for HERMES-{\em relaxed} protocol)
could yield a higher success rate, albeit at the expense of peptide diversity.}

Overall, HERMES demonstrates a potential for designing highly immunogenic peptides up to five amino
acid substitutions away from the wild-type--a task that would otherwise require exploring a vast
sequence space of roughly $10^8$ to $10^{10}$ possibilities for a 9-residue to an 11-residue peptide
(Fig.~\ref{Fig:4}A-C). {This capability raises an important question: to what extent do these
substantial changes in sequence alter the binding geometry of the peptide inside the MHC-TCR cleft?
To this end, we used AF3 to model the TCR-pMHC interfaces for the experimentally validated
immunogenic designs from the three systems. Although we see slight variation in both peptide and TCR
backbone conformations, overall, we found that our designed immunogenic peptides adopt conformations
very close to their wild‑type templates, even when those designs were generated with HERMES‑{\em
relaxed} (Fig.~\ref{Fig:4}E). We attribute this structural convergence to two potential factors: (i)
residual limitations in HERMES peptide‑backbone sampling strategy, and (ii) an inherent AF3 bias
toward conformations it encountered during training. Given that TCRs can cross-react with peptides
of markedly different conformational poses~\cite{Riley2018-zi}, developing methods that can better
explore peptide conformational space within the TCR-MHC groove remains a compelling direction for
future research.}


\subsection*{Structural basis for TCR specificity}
TCRs exhibit substantial degeneracy in their recognition, with some autoimmune T-cells
experimentally shown to recognize over one million distinct peptides~\cite{Wooldridge2012-zq}.
However, the full extent of TCR degeneracy for typical receptors remains unclear, largely due to the
limitations in high-throughput experimental assays for TCR recognition of many peptides presented on
different MHCs.

Our computational peptide design framework helps address this limitation, at least for the TCR-MHC
complexes with known structural templates. Specifically, for a given TCR-MHC pair, we can use
HERMES-{\em fixed} protocol to generate a peptide position weight matrix (PWM), representing the
ensemble of peptides presented by the MHC and recognized by the TCR. We use the entropy of this PWM
as a proxy for TCR degeneracy.

Fig.~\ref{Fig:5}A shows the entropy of these peptide distributions for 105 TCR-MHC pairs across 10
human MHC-I alleles; see Table~\AllPdbsTable~for details on these structures. The median peptide
entropy across these TCR-MHC pairs is approximately 10~bits, indicating that a typical TCR-MHC pair
can recognize on the order of $ \simeq 10^3$ peptides. In contrast, examining the degeneracy of
MHC-I presentation alone in humans--using the peptide PWMs gathered from the MHC Motif
Atlas~\cite{Tadros2023-ij}--yields a median entropy of 31 bits, implying that MHC-I molecules can
typically recognize $\simeq 2\times 10^{9}$ distinct peptides (Fig.~\ref{Fig:5}A). A similar
analysis on 36 TCR-MHC pairs spanning three mouse MHC-I alleles shows comparable trends, with a
MHC-I degeneracy on the order of $\simeq5 \times 10^8$ peptides (29 bits), and the TCR-MHC
degeneracy of $\simeq 10^3$ {peptides} (10 bits); see Fig.~\ref{Fig:5}B, and Table~\AllPdbsTable~for
details on these structures. This pronounced reduction in entropy---observed by comparing the
distribution of peptides recognized by MHC-I molecules alone to that recognized by TCR-MHC-I
complexes---quantifies how TCRs constrain the accessible antigenic shape space, as illustrated by
the polar interactions that they form with peptides in Fig.~\ref{Fig:5}C.

{As an alternative to MHC Motif Atlas~\cite{Tadros2023-ij}, we explored the ability of HERMES-{\em
fixed} to model MHC degeneracy. When structural coverage is extensive (roughly $>30$ structures of
distinct peptides for a given MHC allele), the motifs produced by HERMES align closely with those in
the MHC Motif Atlas (SI, Fig.~\pMHCPWMsFigure, and Table~\pMHCPWMsTInfoable). With sparser
structural coverage, HERMES underestimates the peptide diversity, likely because backbone
conformations are insufficiently sampled. Given the imbalanced distribution of structures across MHC
alleles, we therefore continue to use MHC Motif Atlas~\cite{Tadros2023-ij} as our reference for
computing MHC motif specificity.}

Because our analyses focus on interactions involving typical TCR and MHC molecules, our
entropy-based degeneracy estimates are smaller than those from previous theoretical studies, which
often relied on experimental observations from autoimmune T-cells or limited synthetic peptide
libraries~\cite{Sewell2012-du}. Moreover, the HERMES-{\em fixed} model generates PWMs based on
peptides whose backbone conformations match those of the wild-type template, likely underestimating
the true degeneracy of TCR-MHC complexes (Fig.~\ref{Fig:5}). Although HERMES-{\em relaxed} can
partially address this issue, its designs are only reliable within roughly five amino acids of the
native peptides {with similar backbone conformations (Fig.~\ref{Fig:4})}, preventing a full
exploration of the possible peptide poses for a given TCR-MHC complex. Nonetheless, our estimates
can serve as informative lower bounds on typical MHC-I and TCR degeneracies.

